% \section{Contributin I: Evaluation of Incremental Entity Extraction with Background Knowledge and Entity Linking}
\section{Adapting an Entity Linking Benchmark to Incremental Entity Linking}
\label{sec:ijckg}



The work presented in this section was published in 2022 at the 11th International Joint Conference on Knowledge Graphs (IJCKG 2022)~\cite{ijckg2022}, co-located with the 21st International Semantic Web Conference (ISWC 2022).
%
% The proposed approach addresses the \nilchallenge by extending \ac{el} with a background \ac{kbr} with NIL prediction, novel entity clustering, and the incremental population of the \ac{kr} with novel entities. We apply this methodology to an existing benchmark, propose a baseline pipeline, and demonstrate the increased difficulty of \acl{ee} in incremental settings.
% It focuses on batch-based \acl{inel}, proposing a procedure to adapt existing \ac{nel} datasets to this setting. We applied this method to the WNUM \ac{el} dataset~\cite{Onoe_Durrett_2020} obtaining its incremental version, \emph{Incremental WNUM}, available from GitHub~\footnote{https://github.com/rpo19/Incremental-Entity-Extraction}.
%
At the time of publication (2022), as described in \Cref{sec:incrementalel}, the advancements in the NIL prediction task were left behind the one in \ac{nel}. For instance, the BLINK approach~\cite{blink} was not yet adapted for NIL prediction.
For this reason, addressing the \nilchallenge, we worked towards producing a dataset and an evaluation procedures for aiding the research in this field.

Following this line, during my \phd I co-authored \emph{\mammotab}~\cite{mammotab}, a large-scale dataset of 1 million annotated Wikipedia tables extracted from 20 million pages and aligned with Wikidata. \mammotab is not part of this thesis but still addresses a related topic, \acf{sti}---the field that studies interpretation tasks for tabular data, that are linking cells containing entity mentions to the correct entry in a \ac{kr} (similarly to \ac{nel} for text), identifying the types of table columns, and the relations between different columns of the table~\cite{cremaschi2024surveysemanticinterpretationtabular}. Among the challenges addressed, \mammotab was also constructed for the \nilchallenge, including NIL-labeled instances. This work has been published in the Proceedings of the Fourth edition of the Semantic Web Challenge on Tabular Data to Knowledge Graph Matching (SemTab 2022)~\cite{mammotab}, co-located with the 21st International Semantic Web Conference (ISWC 2022).

Adapting and evaluating existing techniques for \acl{inel} is part of the \ref{ob:extraction_assessment}, which aims at assessing the feasibility of using existing approaches in the legal domain, since the \nilchallenge is relevant for this domain.
The contributions to \acl{inel} are described in the next sections and can be summarized as follows.
\begin{enumerate}
    \item A general methodology to adapt any \ac{nel} dataset to incremental \ac{nel}.
    \item A dataset for \ac{inel}, suitable for evaluating \ac{nel}, NIL prediction, and NIL clustering in the incremental scenario. The dataset was created from an existing \ac{nel} benchmark~\cite{Onoe_Durrett_2020}, applying the proposed adaptation methodology.
    \item A baseline pipeline, constructed on the BLINK bi-encoder~\cite{blink}, additionally supporting NIL prediction and clustering.
    \item Evaluation procedures for incremental \ac{nel}.
    \item An assessment of the performance and main challenges of the baseline pipeline on the incremental scenario.
\end{enumerate}

\todo{dire che non facciamo ner e perche?}

The following sections describe, in order, the baseline pipeline with the motivations behind the choice of each component (\Cref{sec:ijckg:pipeline}), the procedure to adapt a \ac{nel} dataset to the incremental setting and its application to WNUM~\cite{Onoe_Durrett_2020} (\Cref{sec:ijckg:dataset}), followed by an experimental evaluation of the baseline pipeline on the incremental-WNUM and a discussion of the challenges (\Cref{sec:ijckg:experiments}).\todo{divide exp and discussion?}

\subsection{A Baseline Pipeline for Incremental Entity Linking}
\label{sec:ijckg:pipeline}

The proposed baseline pipeline consists of three main modules: \ac{nel}, NIL prediction, and NIL clustering. As illustrated in \Cref{fig:pipeline}, the process operates as follows. Given a mention and its context, the linking module retrieves the most suitable candidate from a \emph{\acf{bgkr}}. The NIL prediction module then verifies whether this candidate is correct. If so, the mention is linked to the \ac{bgkr}; otherwise, it is marked as NIL. All NIL mentions in the current batch are subsequently passed to the NIL clustering module, which groups together those referring to the same unseen entity. Each resulting cluster corresponds to a new entity, which is added to the \emph{\acf{newkr}} to support future linking tasks. Although the distinction between \ac{bgkr} and \ac{newkr} is not strictly necessary---as \ac{inel} could also be achieved by storing all entities in a single reference \ac{kr}---we maintain this conceptual separation for greater clarity.

\begin{figure}[htbp]
    \centering
    \includegraphics[
    width=.35\textwidth,
    keepaspectratio]
    %{imgs/pipeline_simple.drawio.png}
    {images/ijckg/pipeline_simple.png}
    \caption{Schema of the pipeline.
    }
    \label{fig:pipeline}
\end{figure}
\todo{rifare figura con acronimi giusti}

\paragraph{Entity Linking} The \ac{nel} module relies on a bi-encoder architecture~\cite{Humeau2020Poly-encoders,blink,gillick-etal-2019-learning}, which retrieves candidate entities through \ac{ann} search and scores mention--entity pairs using similarity measures such as the dot product.

In our experiments, we used the pre-trained BLINK~\cite{blink} bi-encoder, publicly available on GitHub\footnote{\url{https://github.com/facebookresearch/BLINK}}. For the \ac{bgkr}, we adopted the FAISS~\cite{faiss} index built from the August 2019 Wikipedia dump released by \textcite{blink}, ensuring that entities designated as NIL for the experiments were excluded.

\paragraph{NIL Prediction} Given a mention and the scored entity candidates produced by the linking system, NIL prediction determines whether the top-ranked candidate entity is correct or not. According to this prediction, the mention is either linked to the top-ranked entity or marked as NIL. NIL prediction does not alter the ranking produced by the \ac{nel} module, under the assumption that if the top-ranked candidate according to the \ac{nel} system is incorrect, than the correct candidate is not present in the \ac{kr}.

Our implementation employs a logistic regression model whose input features include the score of the top candidate (the dot product between bi-encoder mention and entity embeddings) and the difference between the top and second-best scores (\textit{secondiff}). This configuration was selected through an ablation study (\Cref{tab:feature_ablation_study}) conducted using the training set for model fitting and the development set for evaluation. The study explored additional features such as textual similarity between the mention and the entity title---using Levenshtein~\cite{levenshtein1966} and Jaccard~\cite{jaccard1901etude} similarities---and basic statistics of the scores of the top-$k$ candidates---namely mean, median, and standard deviation.

The logistic regression estimates $p \in [0,1]$, which represents the estimate probability that the top candidate is correct (values close to~0 indicate NIL). If the prediction is NIL, the mention is passed to the NIL clustering step to identify possible coreferent NIL mentions.
% Note that the bi-encoder scores used in this model are unnormalized (dot-product based) and highly dependent on the embedding space.



\begin{table}
    \centering
\begin{tabular}{lrrr}
\toprule
F$_1$ &  NIL &  $\neg$NIL &  macro F$_1$  \\

\midrule
Max                               &  15.2 &  82.7 &          48.9 \\
Max, Levenshtein, Jaccard           &  23.5 &  82.9 &          53.2 \\
Max, Stats10, Levenshtein, Jaccard   &  50.6 &  83.4 &          67.0 \\
Max, Secondiff, Levenshtein, Jaccard &  51.2 &  83.0 &          67.1 \\
Max, Secondiff                     &  51.4 &  83.1 &          67.3 \\
\bottomrule
\end{tabular}
    \caption{NIL prediction feature ablation study on the dev set (trained on the train set). \textit{max} is the score of the best candidate, \textit{levenshtein} and \textit{jaccard} are textual similarities between the mention and the entity title, \textit{secondiff} is the difference between the best and the second-best score.}
    \label{tab:feature_ablation_study}
\end{table}

\paragraph{NIL Clustering} NIL clustering is performed in multiple stages, considering both the mention surface form and its embedding representation. This approach helps prevent errors due to homonymy. For each cluster of mentions, we then obtain a representation of the corresponding entity using the cluster medoid vector.

We evaluated three clustering strategies. The first, \textit{$\text{GNN}_\text{B}$}, applies a greedy nearest-neighbor (GNN) algorithm on the dense mention vectors obtained from the bi-encoder. Each mention $m$ is clustered with others whose similarity with $m$ exceeds a predefined threshold. This method, combined with various vectorizers, has shown promising results in \textcite{streaming_crossdoc_coreference}.

The second method, \textit{$\text{GNN}_\text{F}$}, also relies on GNN but uses a feature-based vectorizer following \textcite{shrimpton-etal-2015-sampling}. Mentions are encoded as character skip-bigram indicator vectors for surface text and TF-IDF~\cite{jm3} vectors for context. For both $\text{GNN}_\text{B}$ and $\text{GNN}_\text{F}$.

The third method, referred to as \emph{three-step clustering (3-step)}, follows \textcite{simone_24}. First, mentions are grouped by surface similarity: mentions with an edit distance (Damerau-Levenshtein edit distance~\cite{damerau1964}) $\le 3$ are clustered together, except for short mentions of less than 3 characters that are clustered only if identical. This step groups similar mentions but may combine semantically different entities (e.g., \texttt{Lisbon (1956 film)} and \texttt{Lisbon, Capital of Portugal}). Next, we apply single-linkage agglomerative hierarchical clustering~\cite{Rokach2005} within each group based on semantic similarity derived from bi-encoder embeddings, to separate polysemous mentions. Finally, each sub-cluster is represented by its medoid vector, and semantically close sub-clusters are merged based on their medoid similarities, with another single-linkage agglomerative hierarchical clustering pass.

All clustering thresholds were tuned via grid search so that the number of predicted clusters on the development set roughly matched the number of unique entities---as in \textcite{streaming_crossdoc_coreference}.

\paragraph{Novel Entity Representation} New entities lack meaningful textual descriptions, which are typically required by the bi-encoder for representation. However, by leveraging the dense mention embeddings, we can derive a vector representation for a new entity directly from its mentions. In this sense, the entity representation is inferred from real-world usage examples.

We evaluated three strategies for representing new entities from the clusters of NIL mentions, using respectively:
\begin{itemize}
    \item the \textit{first} mention encountered,
    \item the \textit{medoid} vector of all mentions in the cluster at insertion time,
    \item or including \textit{all} mention vectors in the \ac{kb}, treating them as separate entries.
\end{itemize}

\Cref{tab:new_entities_recall} reports the recall values obtained by the bi-encoder for each strategy, together with the corresponding time and storage requirements, as measured on the AIDA dataset~\cite{aida}.
Although including all mention vectors yields the highest recall, we adopt the medoid-based strategy as a compromise between efficiency and effectiveness. It achieves comparable performance while requiring only 5\% of the time and 22\% of the storage of the full approach.
\begin{table}[!htbp]
\centering
\begin{threeparttable}
    \caption{Recall@k results of the linking from scratch experiment (on AIDA~\cite{aida}) with retrieval time and the disk usage of each indexing configuration. The indexes were created using FAISS~\cite{faiss} IndexFlatIP class\tnote{1}.}
    \label{tab:new_entities_recall}
    \begin{tabular}{l|rrrr|rrr}
        \toprule
         & \textbf{R@1} & \textbf{R@3} & \textbf{R@10} & \textbf{R@30} & \textbf{Vectors} & \textbf{Time (s)} & \textbf{Disk Usage (MB)} \\
         \midrule
         First  & 73.7 & 85.6 & 91.9 & 95.8 &  4022 &  2.4 & 16\\
         Medoid & 79.5 & 91.1 & 95.9 & 98.7 &  4022 &  2.4 & 16\\
         All    & 93.9 & 96.7 & 98.5 & 99.3 &  18319 &  44.1 & 72 \\
        \bottomrule
    \end{tabular}
    \begin{tablenotes}
        \footnotesize
        \item[1] \url{https://faiss.ai/cpp_api/struct/structfaiss_1_1IndexFlatIP.html}
    \end{tablenotes}
\end{threeparttable}
\end{table}



% Despite human validation, the challenge of selecting the most representative mentions for new entities remains. Since embeddings are not interpretable, a validator may inadvertently choose non-representative mentions, introducing bias in subsequent batches.

% \subsection{Transforming a \ac{nel} Benchmark into an Incremental Entity Extraction Benchmark}
\subsection{Benchmark Adaptation Procedure: \ac{nel} to \ac{inel}}
\label{sec:ijckg:dataset}
To evaluate our pipeline in a realistic setting, the chosen test dataset should reflect the characteristics expected in real-world scenarios. In particular, it should satisfy the following conditions:

\begin{enumerate}
\item The entity frequency distribution should exhibit a long right tail, where a few entities appear very frequently, while the majority are mentioned rarely.
\item Entities in the training and test sets should belong to the same domain. Indeed, domain adaptation is not in the scope of this section, which focuses on the \nilchallenge.
\item The dataset should be sufficiently large to train data-intensive models and to produce a test set that can be split into multiple batches.
\end{enumerate}

While the first condition may not hold in domain-specific contexts---for instance, in legal or investigative documents, where high-frequency NIL entities unknown to public \acp{kr} may occur---we consider it essential when using general-domain datasets, as in this section. Otherwise, marking high-frequency entities as NIL would be unrealistic. Additionally, approaches based on \acp{llm} may exhibit bias, since these models encode substantial world knowledge (see \Cref{sec:llms_vs_kbs}), which can simplify the \ac{inel} task for entities that appear frequently in their training data.

Among the candidate datasets suitable for adaptation to the incremental scenario are AIDA~\cite{aida}, the Zero-shot EL dataset~\cite{logeswaran-etal-2019-zero}, KORE50~\cite{kore50}, TACKBP-2010~\cite{ji2010overview}, and WikilinksNED Unseen-Mentions (WNUM)~\cite{Onoe_Durrett_2020}. For this work, we focus on WNUM because it is freely available, includes the above characteristics, and links mentions to Wikipedia entities, which aligns with several \ac{nel} systems~\cite{blink,decao2021autoregressive,barba-etal-2022-extend}. AIDA is freely available too, but it is smaller and lacks ground truth for NIL mentions.
%While Wikipedia may not strictly be considered a knowledge base, it is interconnected with Wikidata and DBpedia, allowing information about most Wikipedia entities to be retrieved from proper \acp{kb}.

Importantly, WNUM ensures that no mention-entity pair appears in more than one set (train, dev, or test)~\cite{Onoe_Durrett_2020}, a desirable property for incremental entity evaluation.

\paragraph{New entities.}
To simulate the presence of new entities, we randomly flag certain entities as NIL while preserving the ground truth---necessary for NIL clustering and subsequent linking. To this aim we assign a probability $p_\text{NIL}(e)$ for each entity $e$, depending on the number of times $e$ is mentioned in the training set ($\#e$), and according to $p$ the desired proportion of entities to flag as NIL. $p_\text{NIL}(e)$ is defined as follows.
\begin{equation}
p_{NIL}(e) = p^{\frac{\#e}{M}} \hspace{10pt} \in (0, 1],
\end{equation}
where $M$ is the median entity frequency in the training set. We use the median rather than the mean because the mean is skewed by the long tail of low-frequency entities. This monotonically decreasing function ensures that entities mentioned only once are more likely to be marked as new, reflecting the real-world scenario in which many entities occur rarely, while only a few become popular over a short period.


In this work, we set $p = 0.1$ to obtain 10\% of NIL entities, following~\cite{agarwal_2021_arborescence_arxiv} and apply Bernoulli trial with probability $p_{NIL}(e)$ to decide, for each $e$, whether to flag it as NIL. Entities marked as NIL are removed from the \ac{bgkr} and treated as unknown. To increase the number of new entities in the development and test sets, some mentions of NIL entities are randomly transplanted from the training set. Both development and test sets are assigned 500 new mentions each to ensure robust evaluation (see \Cref{tab:dataset_nils}).
%
Finally, the test set is split into 10 batches using stratified sampling based on entity frequencies and NIL status. \Cref{tab:dataset_stats_extended} reports detailed per-batch statistics.

\begin{table}[htbp]
    \centering
    \caption{Statistics about the dataset before and after the \textit{transplant}.}
    \label{tab:dataset_nils}
    \begin{tabular}{ll|rr|rr}
    \toprule
    && \textbf{Total Mentions} & \textbf{NIL Mentions} & \textbf{Total Entities} & \textbf{New Entities} \\
    \midrule
        \textbf{Before} &Train & 2.2M & (25,744)  & 86,184 & (17,957) \\
        &Dev.   & 10k  & (316)    & 2,397  & (61) \\
        &Test  & 10k  &(307)     & 2514  & (63) \\
        \toprule
       \textbf{After}&Train & 2.008M & (25,365)  & 81,858 & (17,619) \\
       &Dev.   & 100k   & (501)     & 7,105 & (214) \\
       &Test  & 100k   & (501)     & 6,473 & (248) \\
    \end{tabular}
\end{table}
\vspace{-20pt}

\begin{table}[htbp]
    \centering
    \caption{Per-batch statistics about the test set: number of mentions, entities, NIL mentions, new entities, and \emph{Prev~Entities}: new entities already found in a previous batch (that should be linked to an entity previously added to the \ac{newkr})}
    \label{tab:dataset_stats_extended}
\begin{tabular}{l|rrrrrr}
\toprule
{} &  \textbf{Mentions} &  \textbf{Entities} &  N\textbf{IL~Mentions}  &    \textbf{NIL~Entities} &  \textbf{Prev~Entities} \\
\midrule
$b_0$  &          10k &           2.5k &                 50 &                      37 &                                          - &                                          \\
$b_1$  &          10k &           2.5k &                 50 &                      35 &                                         12 &                                          \\
$b_2$  &          10k &           2.5k &                 50 &                      44 &                                         15 &                                          \\
$b_3$  &          10k &           2.5k &                 50 &                      41 &                                         16 &                                          \\
$b_4$  &          10k &           2.5k &                 50 &                      38 &                                         10 &                                          \\
$b_5$  &          10k &           2.5k &                 50 &                      40 &                                         18 &                                          \\
$b_6$  &          10k &           2.5k &                 51 &                      46 &                                         22 &                                          \\
$b_7$  &          10k &           2.5k &                 50 &                      41 &                                         19 &                                          \\
$b_8$  &          10k &           2.5k &                 50 &                      40 &                                         22 &                                          \\
$b_9$  &          10k &           2.5k &                 50 &                      44 &                                         24 &                                          \\
$ALL$ &         100k &           6.5k &                501 &                     248 &                                          - &                                          \\
\bottomrule
\end{tabular}
\end{table}

\begin{figure}[htbp]
    \centering
    \includegraphics[
        width=.8\textwidth,
        keepaspectratio]{images/ijckg/dataset_creation}
    \caption{Construction of the I-\ac{nel} dataset from a \ac{nel} dataset (a): first $p\%$ of mentions from the corpus are flagged as NILs and corresponding entities are removed from the \ac{kb} (b); then observations are just \textit{transplanted} in order to obtain a well-represented distribution for the evaluation (c); finally the test set is split in batches (d).}
    \label{fig:evluation_outcome}
\end{figure}


\subsection{Evaluation Procedure}
\label{sec:ijckg:evaluation_procedure}
\todo{better define metrics ith formulas}

%The evaluation is run in two ways: first on the whole dataset (single-batch) to allow a comparison with other models from literature; then, it is run incrementally on each batch, so that the drop in performance, due to the lack of information provided in the first steps, can be quantified. \newline
%As far as we know, however, there is no other work which uses the same dataset that train the model reducing the available information in any way: other works should be considered only as a reference for the task difficulty.

%In order to test the model in an incremental way, we trained each component of the pipeline, except for the bi-encoder, on the same data (train \& dev sets) and test them on each batch of the test set, presented one by one.
% @Riccardo BLINK has not been trained on unseen mentions. But on its own wikipedia-based dataset. we used a pretrained model.
%The bi-encoder was in fact trained on a dataset based on Wikipedia\cite{blink}: we limit to remove NIL entities from the FAISS index when it retrieves candidates.
%In fact, the updating process, consisting only in the addition of the new vectors in the index, does not require a train phase, and therefore model weights have a constant value during all the test process: the only component which is modified, guaranteeing an incremental procedure, is the index of entities.

% mostrare challenge dell'incremental
To evaluate \ac{inel}, we must consider that the correct handling of a mention $m$ referring to an entity $e$ may depend on previous batches. If $e$ is not present in the \ac{bgkr} at time $t_i$, $m$ should be classified as NIL. However, if a previous batch already contained a mention of $e$, the system should have added it to the \ac{newkr}, and $m$ should be linked to $e$. \Cref{fig:evluation_outcome} summarizes how mentions should be processed depending on the entity they refer to.

The evaluation procedure for the \ac{inel} task should measure both the overall system performance and the performance on subsets of mentions that, in each batch, need to be: (a) linked to the \ac{bgkr}, (b) classified as NIL, or (c) linked to the \ac{newkr} (see \Cref{fig:evluation_outcome}).  
%
Finally, to assess the impact of error propagation across batches, the \ac{inel} evaluation should be comparable to a standard single-batch approach.

\subsubsection{Evaluation Measures} \label{sec:evaluation2}

Evaluating an incremental pipeline is more complex than evaluating a single model, since errors propagate both through the pipeline and across time (i.e., batches).
%
First, we define measures for the overall \ac{inel} pipeline:  

\begin{enumerate}[leftmargin=*, label=(\alph*), ref=(\alph*)]
    \item \textbf{Link to \ac{bgkr}:} accuracy for mentions that should be linked to the \ac{bgkr};
    \item \textbf{NIL:} accuracy for mentions that should be classified as NIL;
    \item \textbf{Link to \ac{newkr}:} accuracy for mentions that should be linked to the \ac{newkr};
    \item \textbf{Overall Accuracy:} accuracy across all mentions.
\end{enumerate}

Secondly, to understand how each component of the pipeline performs on its specific task, we also evaluate the modules separately. In particular, for the \ac{nel} module, we measure Recall@1\todo{dire che in el accuracy e recall@1 coinciding?}. For the NIL predictor, we compute precision, recall, and F$_1$-score for the NIL class. Finally, for NIL clustering, following the co-reference evaluation methodology~\cite{streaming_crossdoc_coreference}, we calculate precision, recall, and F$_1$ for MUC, B$^3$, and CEAF$_e$, also reporting the average F$_1$ score.\todo{spiegare queste metriche? dove? in sota?}


Note that the handling of a mention $m$ referring to an entity $e$ previously added to the \ac{newkr} is considered correct only if it consist linking $m$ to $e$, while labeling $m$ as NIL is considered an error. This ensures that the system does not achieve artificially high scores by treating every mention as a new entity.\todo{rimuovere?}

NIL prediction outcomes can also mitigate \ac{nel} errors. If the linking is incorrect, the NIL predictor should ideally classify the mention as NIL because the suggested entity is wrong---even if the correct entity exists in the \ac{bgkr}\todo{force the nil assumption}. Conversely, if the NIL predictor fails to classify a new entity as NIL, it is considered an error.

Finally, if a false positive in NIL prediction creates a new entity while the correct entity already exists in the \ac{bgkr}, all mentions linked to this incorrect new entity are treated  as errors.

\subsection{Experiments}
\label{sec:ijckg:experiments}

To enable comparison with other models, we tested our baselines using two experiments. The first experiment evaluates the models on the entire unbatched test set in a \emph{single pass} ($\varnothing$), performing only one step of clustering. The second experiment conducts an incremental evaluation over the 10 batches of the test set, which is the main focus of this work and simulates the arrival of new documents to the pipeline. Note that the evaluation procedures differ between the two experiments because NIL clustering adds novel entities to the \ac{newkr} that can be linked in subsequent batches.

The first experiment serves two purposes: it provides results comparable to those reported in the literature, and it allows us to estimate the performance drop introduced by the incremental procedure. We evaluated the three baselines, which differ in their clustering approach, using both experimental setups. In \Cref{tab:superresults}, we report only the NIL clustering performance for all baselines, while the remaining metrics are obtained using the top-performing clustering approach (3-step).

Finally, we conducted an additional experiment, called ``correct'', in which the outputs of previous components are corrected before proceeding through the batches. This setup helps to better analyze error propagation. In \Cref{tab:superresults}, we report the ``correct'' performance only for the top-performing pipeline (3-step).



\begin{figure}
    \centering
    \includegraphics[
    width=.5\textwidth,
    keepaspectratio]
    {images/ijckg/experiments_expected.drawio.png}
    \caption{Schema of the expected outcome of the system, given a mention $m$ referring to an entity $e$, when (a) $e$ is known in the \ac{bgkr}, (b) $e$ has already been added to the \ac{newkr} while processing the previous batches, (c) $e$ is not in \ac{bgkr}, nor in \ac{newkr}.}
    \label{fig:evluation_outcome}
\end{figure}

\subsection{Results}
\label{sec:ijckg:results}
In order to investigate which component of the pipeline has a stronger contribute to the final error, we analyzed each component one by one.
% COMMENTO:
% il link alle nuove entità cala dopo il quinto batch, abbassando il livello totale delle performance
% questo perché le rappresentazioni delle nuove entità sono sub-ottimali rispetto alla distribuzione.
% le performance totali sono migliori, nel senso che non "degenerano" andando avanti con l'aggiungere entità all'indice

%Results are not compared against other models to prevent the temptation of a direct comparison, regardless dataset modifications (as described above) and the different task (which, in our case, is incremental through time).

% Being a binary classification, the NIL predictor performance is measured with precision and recall,
% TODO ENSURE (normalized by the support)
%As shown in \Cref{scores:nil_predictor}.
% The NIL prediction component show high precision scores that grow from $b_0$ to $b_9$. Intuitively, this means that the number of false positives---$\neg$NIL mentions predicted as NIL, in our case---is decreasing over the batches. 

% This happens since the surface form, varies between train and test set~\cite{Onoe_Durrett_2020}
% Errors in the novel entities identification (NIL prediction) and representation (NIL clustering) affect also the \ac{nel} process in the following steps, since the index suffers of ambiguous representations (entities which are linked to a cluster that represents a novel entity which should be linked to a pre-existing one).

The results of the experiments are reported in \Cref{tab:superresults}.
The performance of the \ac{nel} module (BLINK bi-encoder~\cite{blink}) is competitive with more complex models, such as GENRE~\cite{decao2021autoregressive}, which achieves R@1$_\text{total} = 74.7$ and R@1$_\text{unseen} = 70.4$---on mention--entity pairs unseen at training time---values that are very close to ours (excluding error propagation).

As expected, the observed performance drop is mainly due to error propagation in the incremental procedure rather than to limitations of the \ac{nel} model itself. 
Indeed, in the ``correct'' experiment, where errors are prevented from propagating, the performance remains stable across batches.

The NIL prediction component, despite its high precision, suffers from low recall, leading to a large number of false negatives---or \emph{false links}, that is, mentions incorrectly predicted as not NIL ($\neg$NIL) and linked to the \ac{kr} while actually referring to new entities.
Errors from the NIL predictor introduce two critical issues in the pipeline.
First, each false positive---or \emph{false NIL}, i.e., a $\neg$NIL mention incorrectly classified as NIL---creates a spurious entity.
This results in redundant representations of the same entity in the \ac{newkr} and increases error propagation across batches, since mentions linked to these spurious entities in later steps are counted as errors.
For this reason, to mitigate error propagation over time, achieving higher precision is preferable to higher recall, as \emph{false links} count as local errors but do not cause propagation by polluting the \ac{newkr} with redundant entities.

\Cref{fig:clusters} shows the number of predicted clusters as a function of their size, together with the expected distribution. 
The general agreement between the two indicates that the pipeline does not overproduce clusters, and that the high precision yields cluster sizes broadly consistent with the expected values.



% Regarding NIL clustering, analyzing and understanding this component is crucial, since it produces side effects that propagate through subsequent steps, leading to different outcomes when evaluated incrementally versus in a single pass.
Among the evaluated approaches, the top-performing one is the \textit{3-step} method, likely because it combines both lexical and semantic similarities between mentions, while the other two ($\text{GNN}_\text{B}$ and $\text{GNN}_\text{F}$) rely exclusively on dense semantic representations.


% Results, as expected, are higher in the single-pass experiment than in the incremental one: this is due to the propagation of error through time in the index updating.
% In fact, the high B3 Recall is due to the high frequencies of single-mentioned entities in the dataset which are presented to the NIL clustering component (remember that the function that assigns NIL values is a decreasing monotonous function).
% This translates in a high number of clusters correctly composed by a single mention.
% Having a suboptimal representation of those entities, which is given by a single observation and not by a textual description or more examples, the index begins suffering due to data quality problems.



\begin{figure}
    \centering
    \includegraphics[width=0.8\linewidth]{images/ijckg/clusters.png}
    \caption{The absolute frequency of clusters by size (the points) vs the expected values (the curve) from the test-set.}
    \label{fig:clusters}
\end{figure}



\begin{landscape}
%%%%%%%%%%%%%%% NEWWW
% incremental_3step
\begin{table*}
    \centering
\begin{tabular}{llrrrrrrrrrrrr}
\toprule
{} &   &  $b_0$  &     $b_1$  &     $b_2$  &     $b_3$  &     $b_4$  &     $b_5$  &     $b_6$  &     $b_7$  &     $b_8$  &     $b_9$  &     $\overline{b}$ & $\varnothing$ \\
\midrule
% Three step    &&&&&&&&&&&& \\
% \hline
% & NEL &&&&&&&&&&&& \\
\textbf{NEL} & R@1          &  72.64 &  62.03 &  55.43 &  51.68 &  46.53 &  45.02 &  41.71 &  40.56 &  39.03 &  38.39 &  49.30 & 72.91 \\
\midrule
% & NIL Pred &&&&&&&&&&&& \\
\textbf{NIL} & P            &  65.16 &  66.25 &  70.55 &  74.30 &  78.92 &  78.64 &  80.00 &  79.53 &  81.38 &  81.96 &  75.67 & 64.01 \\
\textbf{Pred} & R            &  46.74 &  30.51 &  30.55 &  31.47 &  33.59 &  34.14 &  34.32 &  34.21 &  35.79 &  36.10 &  34.74 & 46.06 \\
& F$_1$           &  54.43 &  41.78 &  42.63 &  44.22 &  47.13 &  47.61 &  48.04 &  47.84 &  49.71 &  50.12 &  47.35 & 53.57 \\
\midrule
% & NIL Clust    &&&&&&&&&&&& \\
% & Three step    &&&&&&&&&&&& \\
\multicolumn{14}{l}{\textbf{NIL Clust:} comparison among 3-step, $\text{GNN}_\text{B}$, $\text{GNN}_\text{F}$: best results highlighted in bold}  \\
3-step & MUC F$_1$       &  \textbf{96.84} &  \textbf{95.97} &  \textbf{95.89} &  \textbf{96.49} &  \textbf{96.96} &  \textbf{97.10} &  \textbf{96.94} &  \textbf{97.48} &  \textbf{97.26} &  \textbf{97.48} &  \textbf{96.84} & \textbf{97.38} \\
& B3 F$_1$        &  \textbf{97.43} &  \textbf{97.07} & \textbf{96.73} &  \textbf{96.25} &  \textbf{96.48} &  \textbf{96.66} &  \textbf{96.11} &  \textbf{97.14} &  \textbf{95.79} &  \textbf{96.55} &  \textbf{96.62} & \textbf{94.36} \\
& CEAF F$_1$      &  \textbf{95.16} &  \textbf{94.43} &  \textbf{93.10} &  \textbf{93.02} &  \textbf{92.56} &  \textbf{93.15} &  \textbf{92.38} &  \textbf{93.35} &  \textbf{92.49} &  \textbf{93.26} &  \textbf{93.29} & \textbf{82.27} \\
& macro F$_1$       &  \textbf{96.48} &  \textbf{95.82} &  \textbf{95.24} &  \textbf{95.25} &  \textbf{95.33} &  \textbf{95.63} &  \textbf{95.14} &  \textbf{95.99} &  \textbf{95.18} &  \textbf{95.76} &  \textbf{95.58} & \textbf{91.34} \\
\hline
% & Greedy    &&&&&&&&&&&& \\
$\text{GNN}_\text{B}$ & MUC F$_1$       &  86.48 &  84.22 &  84.55 &  87.65 &  89.23 &  89.60 &  89.30 &  91.07 &  90.88 &  91.05 &  88.40 & 91.56 \\
& B3 F$_1$        &  91.33 &  90.12 &  89.07 &  90.34 &  90.92 &  90.74 &  90.36 &  91.69 &  91.72 &  91.92 &  90.82 & 86.46 \\
& CEAF F$_1$      &  84.96 &  83.94 &  80.62 &  81.65 &  80.95 &  81.58 &  80.77 &  82.57 &  81.38 &  81.14 & 81.96  & 64.61  \\
& macro F$_1$ &  87.59 &  86.09 &  84.75 &  86.54 &  87.03 &  87.31 &  86.81 &  88.44 &  87.99 &  88.04 &  87.06 & 80.88 \\
\hline
% & Feature    &&&&&&&&&&&& \\
$\text{GNN}_\text{F}$ & MUC F$_1$       &  36.08 &  31.05 &  31.04 &  29.84 &  28.37 &  28.16 &  27.67 &  28.53 &  26.69 &  29.17 &  29.66 &  65.18 \\
& B3 F$_1$        &  71.66 &  65.16 &  58.68 &  56.80 &  55.03 &  54.78 &  53.31 &  54.40 &  51.89 &  51.16 &  57.29 & 48.83 \\
& CEAF F$_1$      &  62.72 &  56.91 &  49.09 &  47.01 &  44.90 &  44.60 &  43.24 &  44.22 &  41.66 &  40.46 &  47.48 & 35.89 \\
& macro F$_1$ &  56.82 &  51.04 &  46.27 &  44.55 &  42.77 &  42.51 &  41.41 &  42.39 &  40.08 &  40.26 &  44.81 & 49.96 \\
\hline
\textbf{Overall} & (a) Link     &  65.67 &  56.05 &  49.68 &  46.36 &  41.69 &  39.87 &  36.67 &  35.29 &  33.99 &  33.46 &  43.87 & \\
& (b) NIL      &  42.00 &  38.46 &  51.35 &  32.35 &  60.00 &  50.00 &  41.18 &  46.67 &  44.44 &  37.93 &  44.44 & \\
& (c) Link$_{\text{New}}$ &   n/a* &  36.36 &  76.92 &  68.75 &  73.33 &  70.00 &  70.59 &  85.00 &  69.57 &  61.90 &  61.24 & \\
& (d) Acc &  65.55 &  55.96 &  49.72 &  46.35 &  41.80 &  39.96 &  36.74 &  35.42 &  34.10 &  33.53 &  43.91 & \\
\midrule
\hline
\multicolumn{14}{l}{\textbf{Correct}} \\
\textbf{NEL} & R@1          &   72.64 &  72.11 &   72.24 &   72.79 &  72.26 &  73.13 &   71.71 &   72.77 &  71.87 &   72.28 &  72.38 & \\
\hline
\textbf{NIL} & P            &   55.43 &  55.17 &   54.87 &   54.26 &  55.52 &  55.61 &   55.96 &   56.21 &  55.83 &   54.38 &  55.32 & \\
\textbf{Pred} & R            &   43.32 &  42.43 &   45.47 &   43.02 &  43.57 &  43.44 &   42.22 &   44.91 &  42.30 &   43.08 &  43.38 & \\
& F$_1$           &   48.63 &  47.97 &   49.73 &   47.99 &  48.83 &  48.78 &   48.13 &   49.93 &  48.13 &   48.08 &  48.62 & \\
\hline
\textbf{NIL} & MUC F$_1$       &  100.0 &  80.00 &    -- &  100.0 &  66.67 &  66.67 &    -- &  100.0 &   -- &    -- &  51.33 & \\
\textbf{Clust} & B3 F$_1$        &  100.0 &  97.26 &  100.0 &  100.0 &  95.65 &  97.78 &  100.0 &  100.0 &  97.14 &  100.0 &  98.78 & \\
3-step & CEAF F$_1$      &  100.0 &  96.89 &  100.0 &  100.0 &  94.53 &  96.12 &  100.0 &  100.0 &  95.24 &  100.0 &  98.28 & \\
& macro F$_1$       &  100.0 &  91.38 &   66.67 &  100.0 &  85.62 &  86.86 &   66.67 &  100.0 &  64.13 &   66.67 &  82.80 & \\
\bottomrule
\end{tabular}
    \caption{Evaluation results: \textit{\ac{el}}, \textit{NIL Pred}, \textit{Overall}, and \textit{Correct} are obtained using the pipeline with the 3-step clustering algorithm (the top performing one). \textit{Correct} results are obtained correcting the output of the previous components. $\varnothing$ refers to the unbatched single-pass experiment. NIL prediction performance are calculated considering correct when a NEL error is mitigated. *Nothing can be linked to previously added entities in $b_{0}$. ``--'' represents cases where the gold standard contains one element per cluster; in this case, \textit{MUC F$_1$} score (= 0) is not meaningful~\cite{recasens_hovy_2011}.}
    \label{tab:superresults}
\end{table*}
\end{landscape}

These errors can be mitigated through a \acl{hitl} validation process, where a user, using dedicated \acp{ui}, can merge clusters referring to the same entity or split those that incorrectly group different mentions.\todo{This functionality has been implemented in our system.}



\subsection{Discussion: Challenges in Incremental Entity Extraction} \label{sec:futures}
% hitl


%Results obtained in this work show that the approach is promising, using pre-trained neural networks, but can be further improved by the addition of new modules and improvements of existing ones.

%In particular, as already described above, a transformer-based model able to distinguish between mention-mention could significantly improve performances of the whole pipeline, with both a re-ranking of the candidates and providing a better distance score during NIL clustering.
% In addition, it can be used also as a distance score between mentions in order to cluster them in the NIL clustering step.
%This module, however, increases the whole complexity and computational cost at inference time, losing the interactivity of the pipeline.


% even if the classifier obtains relatively high precision levels (see \Cref{tab:superresults}) the error rate is high enough to degrade the performance in the subsequent batches.

%In addition, the NIL predictor, which at the moment uses a logistic regression, may be improved by the use of linear transformations, since the use of non-linear and more complex models (such as random forest) offer no better results.
%A future version may include more features, based on candidates distribution, as well as a standardization layer so that it can be pre-trained on a different dataset.
%All those improvements however are not trivial both to train and test in order to obtain a model not based on dataset scores.

% \begin{enumerate}
%     \item Improve NIL prediction.
%     \item Improve NIL clustering.
%     \item Further investigate on how the system behaves with long-tailed entities.
% \end{enumerate}

\label{sec:ijckg:discussion}

The results show that error propagation across batches is a major challenge in incremental entity extraction. A performance drop is especially evident in the \ac{nel} results. As reported in \Cref{tab:superresults}, the first batch achieves scores comparable to those of the single-pass ($\varnothing$) and ``correct'' experiments, but performance progressively degrades in subsequent batches.

As discussed earlier, a main source of errors lies in false positives (false NILs) introduced during NIL management. This confirms that accurate NIL prediction is a crucial challenge in the \ac{inel} task: lower precision amplifies the propagation of errors throughout the pipeline, ultimately degrading performance.


In our experiments, novel entities are represented by a single vector---the medoid of their NIL mention cluster. However, evidence from \Cref{tab:new_entities_recall} indicates that this representation is sub-optimal, although more efficient. Exploring alternative strategies, such as using multiple vectors per cluster constitutes a promising direction for future work to achieve a better trade-off between effectiveness and efficiency.

Another open problem concerns when to update the \ac{kr}. In the current setup, the index is updated after processing the first batch in which a new mention appears, but collecting additional observations might yield more reliable representations, possibly guided by a confidence score.
% We leave this aspect for future work, as it is tightly connected to the long-tail entity problem (i.e., entities mentioned only once). In fact, our dataset was enriched with many such entities.
Balancing the need for informative representations with the challenge of handling rare entities remains an interesting avenue for further research.


\subsection{Conclusion}
\todo{remove and put conclusion at the end of chapter}
In this work, we introduce the task of incremental \ac{nel} (\ac{inel}), providing both a dataset and a simple pipeline, both as baseline and as a workbench to identify critical components.
%&The purpose of the task is to reduce the human effort in updating existing \ac{kb}s, automatically obtaining a representation of novel entities.

% The purpose of the pipeline is not only linking mention to entities but also represent novel entities in the same vector space.
Our experiments show that a main challenge for \ac{inel} is to deal with the incremental propagation of error, which is due, especially, to the difficulty of linking entity mentions found in one batch to novel entities, i.e., entities not in the \ac{bgkr} and added to the \ac{newkr} in previous batches. 
Finding representations of novel entities that provide a reasonable trade-off between efficiency and effectiveness is not trivial.
%A sub-optimal representation that may be fast in retrieval, may also introduce errors that propagate incrementally.
%, obtaining worse and worse performance.
% We presented our work as a strong baseline for the task of Incremental \ac{nel} (I\ac{nel}), providing a modification of a ``static'' dataset~\cite{Onoe_Durrett_2020} with a canonical order for data ingestion.
Indeed, the performance of the \ac{nel} component~\cite{blink} is reasonably good at the beginning but worsen over time. Another component that must be improved to deliver robust incremental \ac{nel} is NIL prediction.      
%used in our baseline pipeline has good initial performance but , behaves worse due to error propagation, showing the difficulty of the tasks of novel entities recognition and representation.

\todo{future works removed}
% Future works include the application of our methodology to generate incremental versions of more datasets. In addition, we plan to elaborate on methods to address the main challenges found in I-\ac{nel}, especially in relation to the management of novel entities (NIL prediction and NIL mention representation).  
% %, and deeper studies to improve novel entities representation.
% % and additional studies on when to add a new entity to the KB.

% The dataset, the source code of the procedure to create it, and the code for the evaluation are openly available at \url{https://github.com/rpo19/Incremental-Entity-Extraction/}.